% Results (P8.7)
\section{Results}\label{sec:results}

\subsection{Land classification}\label{sec:res-classification}

\paragraph{Choosing the clusters.} K-Means and Gaussian mixture models were run for $k=3$ to 10
on seven standardised inputs: the cell's shares of water, trees, crops, built and rangeland, its
mean elevation and its mean slope (Figure~\ref{fig:kscan}). K-Means separates the cells far
better at every $k$: silhouette 0.43--0.61 \citep{rousseeuw1987}, against 0.26--0.36 for the
mixture models, which also have a worse Davies--Bouldin index \citep{davies1979} at every $k$. The
silhouette of K-Means is flat from $k=5$ to $k=8$ and drops sharply at $k=9$. We use
\textbf{$\boldsymbol{k=7}$} (silhouette 0.606 on the 2023 features, 0.585 on 2018), which is the
best or within 0.003 of the best in both years. The seven clusters are the five land-cover groups
plus two sub-groups that appear in both years:
\begin{itemize}
  \item \textbf{steep natural vegetation} (mean slope \ang{13.2}, 3\,\% of the cells): the
        Bannerghatta hills, mostly ESRI rangeland;
  \item \textbf{the low western valley} (mean elevation \SI{759}{\m}, 3\,\% of the cells),
        where cropland and built land are mixed.
\end{itemize}

\begin{figure}[ht]
  \centering
  \includegraphics[width=\textwidth]{chart_k_scan}
  \caption{Choosing the number of clusters (2023 features). The dashed line marks $k=7$.}
  \label{fig:kscan}
\end{figure}

\paragraph{From clusters to classes.} Each cluster is labelled from its centroid
(Figure~\ref{fig:profiles}): water-dominated clusters are excluded; built-up where built land is
both above 30\,\% and the largest land group; forest where trees exceed 50\,\%, or where trees
and rangeland together exceed 50\,\% on slopes above \ang{8}; everything else is usable, with
cropland-dominated clusters flagged as agricultural. The second forest rule is what keeps the
Bannerghatta scrub forest out of the usable class. Finally, the exclusion mask (water, wetland,
slopes over \ang{15}, the national park) is applied cell by cell.

\begin{figure}[ht]
  \centering
  \includegraphics[width=0.95\textwidth]{chart_cluster_profiles}
  \caption{Cluster centroids (2023): mean shares, elevation and slope, coloured by z-score.
           Each row is named with the class its rule gave and its share of the cells.}
  \label{fig:profiles}
\end{figure}

\paragraph{The three-class map.} Figure~\ref{fig:class3} shows the result for 2023. Built-up land
covers \SI{348}{\km\squared}, usable land \SI{138}{\km\squared}, forest \SI{29}{\km\squared}, and
\SI{66}{\km\squared} are excluded. Between the 2018 and 2023 maps, built-up land grows by
\SI{42}{\km\squared} and usable land shrinks by \SI{44}{\km\squared}
(Figure~\ref{fig:class-areas}).

Against ESRI's own classes, collapsed to the same three, the map agrees on 96.6\,\% of the land
cells. This is expected, since the ESRI fractions are inputs, and it mainly shows that the
cluster rules don't distort them. Against the independent WorldCover map the agreement is
64.0\,\%. Almost all of the difference comes from the products themselves: WorldCover calls only
33\,\% of the area built-up, against 62\,\% in ESRI.

\begin{figure}[ht]
  \centering
  \includegraphics[width=0.7\textwidth]{map_class3_final}
  \caption{Three-class land map, 2023 (final run).}
  \label{fig:class3}
\end{figure}

\begin{figure}[ht]
  \centering
  \includegraphics[width=0.65\textwidth]{chart_class_areas}
  \caption{Class areas in the 2018 (validation) and 2023 (final) maps.}
  \label{fig:class-areas}
\end{figure}

\subsection{Building similarity}\label{sec:res-similarity}

\paragraph{Single-building queries.} Figure~\ref{fig:queries} shows two of the four example
queries from notebook~05, for a school and an industrial building. The most similar usable cells
lie at the edges of the built-up area, often several kilometres from A, wherever the mix of
built-up land, roads and terrain in the 250 and \SI{500}{\m} rings matches A's. Each result
comes with an explanation: for the industrial building's best match, the distances to built-up
land, water and a major road differ by less than 0.1 standard deviations, while the built-up
share within \SI{250}{\m} differs most (0.6 standard deviations: 80\,\% instead of 100\,\%). Land-cover shares that are zero at
both cells are listed last, since such matches say nothing. Euclidean and cosine distance share
14--18 of the top 20 cells; cosine ignores how extreme a profile is, so it is kept only as a
comparison.

\begin{figure}[ht]
  \centering
  \begin{subfigure}[t]{0.49\textwidth}
    \includegraphics[width=\textwidth]{map_similarity_query_school_1019}
  \end{subfigure}\hfill
  \begin{subfigure}[t]{0.49\textwidth}
    \includegraphics[width=\textwidth]{map_similarity_query_industrial_274244}
  \end{subfigure}
  \caption{Building similarity queries: the 20 usable cells whose surroundings are most similar
           to a school (left) and an industrial building (right), A marked with a star.}
  \label{fig:queries}
\end{figure}

\paragraph{Sanity checks.} For 200 random buildings, a cell within one cell of the building is in
the top 1\,\% of all cells even when the building's own cell is left out: surroundings vary
smoothly in space, so the profile identifies its location. Reference cells held out by
\SI{2}{\km} block are ranked above the other cells with AUC 0.895 by the nearest-neighbour
similarity, 0.984 by a random forest and 0.887 by distance to built-up alone. Keeping both ring
scales beats the \SI{500}{\m} rings alone (0.895 vs 0.871).

\paragraph{Aggregate similarity.} Averaged over the ten nearest reference profiles, the similarity
of every candidate cell reaches \textbf{AUC 0.773} on the growth by 2022/23
(Table~\ref{tab:auc}), above distance to built-up land, although it never sees a growth label.

\subsection{Learned models}\label{sec:res-rf}

The \textbf{change RF} reaches \textbf{AUC 0.779}, slightly above distance to built-up land
(Table~\ref{tab:auc}). Its most important input is the built-up share within \SI{500}{\m},
followed by trees and crops nearby and the distance to built-up land; terrain and water add
little (Figure~\ref{fig:importance}). Scored on its own training cells, it would claim AUC
0.997: most of its 2020/21 positives also grew by 2022/23, so in-sample scores would leak the
answer, and out-of-fold scoring is essential. The \textbf{state RF} separates built from
non-built cells almost perfectly (spatial-CV AUC 0.998), but that skill does not transfer to
predicting the next growth (0.702).

\begin{figure}[ht]
  \centering
  \includegraphics[width=0.85\textwidth]{chart_rf_importance}
  \caption{What the change RF relies on: permutation importance (AUC drop on held-out folds) and
           impurity importance.}
  \label{fig:importance}
\end{figure}

\subsection{Combined score, classes and candidate sites}\label{sec:res-suitability}

The weight chosen on the 2020/21 training label is 0.3 for the similarity and 0.7 for the change
RF. In the final run, the 13{,}814 usable cells fall into 1{,}382 S1, 2{,}762 S2, 4{,}144 S3 and
5{,}526 N cells (Figure~\ref{fig:suitability}). Highly suitable land concentrates on the southern
and eastern fringe of the built-up area, along the Hosur Road corridor, and in the low western
valley.

Merging edge-adjacent S1 cells gives 557 candidate patches covering 1{,}382\,ha; 253 are mostly
cropland and flagged as agricultural. Half of the patches are single \SI{1}{\hectare} cells,
often pockets inside built-up areas, so the top sites are taken from patches of at least
\SI{5}{\hectare}. The 20 top sites (Table~\ref{tab:sites}, 138\,ha) are all gently sloping
(\ang{1}--\ang{4}), directly next to existing built-up land, and covered by rangeland or crops.
None of them overlaps an OSM lake (0.06\,ha in total) or the national park.

\begin{figure}[ht]
  \centering
  \includegraphics[width=0.72\textwidth]{map_suitability_classes}
  \caption{Suitability classes in the final run (2023 data) with the outlines of the candidate
           patches. Forest, built-up and excluded land is always N.}
  \label{fig:suitability}
\end{figure}

\begin{table}[ht]
  \centering\small
  \caption{The ten best candidate sites (patches of at least \SI{5}{\hectare}, ranked by mean
           score). Distances are means over the patch's cells.}
  \label{tab:sites}
  \begin{tabular}{@{}rrrrrrll@{}}
    \toprule
    Rank & Area (ha) & Score & Slope (\textdegree) & Major road (m) & Built-up (m) & Cover & Location \\
    \midrule
    1 & 5 & 0.98 & 1.9 & 371 & 56 & rangeland & 77.64\,E, 12.86\,N \\
    2 & 6 & 0.98 & 1.6 & 1{,}621 & 55 & rangeland & 77.69\,E, 12.78\,N \\
    3 & 5 & 0.97 & 1.9 & 226 & 55 & crops & 77.73\,E, 12.86\,N \\
    4 & 5 & 0.97 & 3.9 & 572 & 57 & rangeland & 77.53\,E, 12.90\,N \\
    5 & 6 & 0.96 & 1.8 & 358 & 55 & rangeland & 77.71\,E, 12.76\,N \\
    6 & 6 & 0.96 & 1.6 & 161 & 55 & rangeland & 77.70\,E, 12.86\,N \\
    7 & 7 & 0.96 & 1.5 & 503 & 55 & rangeland & 77.66\,E, 12.76\,N \\
    8 & 6 & 0.96 & 1.5 & 1{,}265 & 55 & crops & 77.67\,E, 12.79\,N \\
    9 & 6 & 0.96 & 1.4 & 331 & 55 & rangeland & 77.71\,E, 12.74\,N \\
    10 & 5 & 0.95 & 1.3 & 412 & 61 & rangeland & 77.73\,E, 12.79\,N \\
    \bottomrule
  \end{tabular}
\end{table}

\subsection{Validation of all scores}\label{sec:res-validation}

Table~\ref{tab:auc} compares every score on the same 12{,}777 candidate cells, of which 1{,}125
were built in both 2022 and 2023. The \textbf{combined score} is the best: AUC 0.797, and its top
10\,\% of cells contain 3.45 times as much growth as a random tenth. Both PRD targets are met
(best model at least as good as distance to built-up; lift of at least 2 in the top 10\,\%). The
S1 and S2 classes of the validation run hold 65\,\% of the growth on 25\,\% of the evaluated
cells.

\begin{table}[ht]
  \centering\small
  \caption{Validation run: every score computed with 2018 data, tested on persistent growth by
           2022/23 (12{,}777 candidate cells, 1{,}125 grew). RF scores are out of fold.}
  \label{tab:auc}
  \begin{tabular}{@{}lccccc@{}}
    \toprule
    Score & AUC & Lift top 10\,\% & Lift top 30\,\% & AUC adjacent & AUC outlying \\
    \midrule
    Combined score & \textbf{0.797} & \textbf{3.45} & \textbf{2.40} & \textbf{0.822} & 0.512 \\
    Change RF & 0.779 & 3.23 & 2.27 & 0.803 & 0.512 \\
    Building similarity & 0.773 & 2.99 & 2.22 & 0.796 & 0.511 \\
    Distance to built-up (baseline) & 0.763 & 3.04 & 2.15 & 0.792 & 0.439 \\
    State RF & 0.702 & 3.29 & 1.92 & 0.722 & 0.484 \\
    MCDA / AHP & 0.636 & 1.86 & 1.64 & 0.655 & 0.420 \\
    MCDA without land cover & 0.614 & 1.73 & 1.46 & 0.630 & 0.430 \\
    Random & 0.499 & 1.00 & 1.02 & 0.501 & 0.477 \\
    \bottomrule
  \end{tabular}
\end{table}

\begin{figure}[ht]
  \centering
  \begin{subfigure}[t]{0.49\textwidth}
    \includegraphics[width=\textwidth]{chart_roc}
    \caption{ROC curves.}\label{fig:roc}
  \end{subfigure}\hfill
  \begin{subfigure}[t]{0.49\textwidth}
    \includegraphics[width=\textwidth]{chart_toc}
    \caption{Total operating characteristic.}\label{fig:toc}
  \end{subfigure}
  \caption{Predicting growth from 2018.}
\end{figure}

\paragraph{Growth type.} Every score finds growth next to existing built-up land well (AUC
0.79--0.82 for the best ones) but growth away from it only at chance level (0.51); distance to
built-up and MCDA even fall below chance there (Figure~\ref{fig:growth-type}). This matches the
finding that diffusional growth is much harder to predict than aggregated growth
\citep{wang2021}. Only 91 of the 1{,}125 grown cells are outlying, so this hardly moves the
overall AUC.

\paragraph{Ablation.} Dropping the built-up context (built-up shares in the rings and distance to
built-up) costs the most: $-0.034$ AUC for the change RF and $-0.057$ for the similarity
(Figure~\ref{fig:ablation}). Roads and water add nothing once the built-up context is known.

\begin{figure}[ht]
  \centering
  \begin{subfigure}[t]{0.52\textwidth}
    \includegraphics[width=\textwidth]{chart_growth_type}
    \caption{AUC by growth type.}\label{fig:growth-type}
  \end{subfigure}\hfill
  \begin{subfigure}[t]{0.46\textwidth}
    \includegraphics[width=\textwidth]{chart_ablation}
    \caption{AUC change when a feature group is dropped.}\label{fig:ablation}
  \end{subfigure}
  \caption{Where the models succeed and what they rely on.}
\end{figure}

\subsection{MCDA / AHP baseline}\label{sec:res-mcda}

The classic method uses five criteria, each rescaled to 0--1: distance to existing built-up land,
distance to a major road, slope, land cover (rangeland and bare ground easiest to develop, then
crops, then trees) and a setback from water. The AHP pairwise matrix gives built-up and road
distance 0.30 each, slope and land cover 0.16 each and water 0.09. Its consistency ratio of 0.003
is far below Saaty's 0.1 limit \citep{saaty1980}.

Tested against persistent growth from 2018 to 2022/23, the overlay reaches \textbf{AUC 0.636}
(Table~\ref{tab:auc}), with a lift of only 1.86 in the top 10\,\%, below a single criterion, distance to built-up land (0.763). Scored one at
a time, only that criterion carries real signal: slope reaches 0.612, and roads, land cover and
water stay between 0.52 and 0.55, close to chance. The AHP weights give the weak criteria 70\,\%
of the weight, which dilutes the one strong signal. We deliberately did not tune the weights on
the growth label, which is the test set. A fixed, defensible expert overlay is exactly what this
baseline stands for, and it illustrates the weakness the literature describes: the weights decide
the map, and nothing checks them against outcomes \citep{malczewski2004}.

\subsection{Sensitivity}\label{sec:res-sensitivity}

Three settings were varied (Figure~\ref{fig:sensitivity}):
\begin{itemize}
  \item \textbf{Slope threshold} (\ang{10}, \ang{15}, \ang{20}): the number of slope-excluded
        cells changes from 1{,}685 to 158, but every AUC moves by less than 0.01.
  \item \textbf{Number of clusters} ($k=5$ to 8): class shares change by about one percentage
        point and agreement with ESRI stays at 96.6--98.8\,\%. $k=5$ agrees slightly better
        with ESRI but loses the steep-vegetation cluster, so forest falls from 5.0\,\% to
        3.9\,\% of the area.
  \item \textbf{Cell size} (50, 100, \SI{200}{\m}; features and labels rebuilt each time): AUCs
        are not comparable across sizes, because the growth label itself changes with the cell,
        and at \SI{200}{\m} only 148 cells grew.
\end{itemize}
What holds in every setting is the ranking: \textbf{change RF $>$ distance to built-up $>$
MCDA}.

\begin{figure}[ht]
  \centering
  \includegraphics[width=0.95\textwidth]{chart_sensitivity}
  \caption{Temporal AUC by slope threshold and by cell size.}
  \label{fig:sensitivity}
\end{figure}
